% !TeX program = xelatex
\documentclass[12pt]{nextart_zh}

\UseTemplateSet{
  layout = zh_doc,
  globalfonts = {cmu,shanggu},
  features = {tables,image}
}

\title{中文論文模板示例}
\author{作者姓名}
\date{\today}

\begin{document}

\maketitle

\begin{abstract}
本文示範 IMPE LaTeX System 中文論文模板的基本起稿方式。它提供標題、摘要、章節結構、
表格與圖片佔位等常見元素，方便後續直接改寫成實際文稿。
\end{abstract}

\section{研究背景}

這一節可用來簡述研究背景、問題意識與材料來源。若需要中外文混排，也可以直接
在正文中混用 English terms、專有名詞與引文說明。

\section{研究問題}

本文可圍繞以下幾點展開：

\begin{enumerate}
  \item 研究對象的文本特徵與語言背景。
  \item 現有研究的整理與不足。
  \item 本文希望提出的方法與觀察。
\end{enumerate}

\section{示例表格}

\begin{table}[htbp]
  \centering
  \begin{tabular}{lll}
    \toprule
    項目 & 說明 & 備註 \\
    \midrule
    材料 & 文本樣本 & 可替換 \\
    方法 & 描寫與比較 & 可擴展 \\
    結果 & 初步觀察 & 可細化 \\
    \bottomrule
  \end{tabular}
  \caption{基本表格示例}
\end{table}

\section{示例圖片}

\begin{figure}[htbp]
  \centering
  \rule{0.7\linewidth}{0.28\linewidth}
  \caption{圖片佔位示例}
\end{figure}

\section{結語}

這份模板的目的是提供一個可以直接開始撰寫的中文論文骨架，而不是只展示接口。
後續可在此基礎上加入引文、腳註、參考文獻與更多 feature。

\end{document}
